%% Đề 48 -- de-tuyen-sinh-lop-10-thpt-mon-toan-nam-2022-2023-so-gddt-lang-son.pdf
\deheader{ĐỀ SỐ 48}%
  {Kỳ thi Tuyển sinh vào lớp 10 THPT — Năm học 2022--2023}%
  {Sở GD\&ĐT Lạng Sơn | Chỉ phần Đại số}
\addcontentsline{toc}{section}{Đề 48 — Tuyển sinh vào lớp 10 THPT — Năm học 2022--2023}

\begin{de_cau}{Câu 1}{2,5 điểm}
\begin{enumerate}[label=\alph*)]
\item Tính giá trị của các biểu thức:
\[A = \sqrt{81} - \sqrt{16};\qquad B = \sqrt{(\sqrt{11}+2)^2} - \sqrt{11}.\]
\item Cho biểu thức: $P = \left(\dfrac{1}{\sqrt{a}-1} + \dfrac{2}{a-\sqrt{a}}\right) \cdot \dfrac{a-1}{\sqrt{a}+2}$ với $a > 0,\; a \neq 1$.
\begin{enumerate}[label=\arabic*.]
\item Rút gọn biểu thức $P$.
\item Tính giá trị của $P$ khi $a = 3 + 2\sqrt{2}$.
\end{enumerate}
\end{enumerate}
\end{de_cau}

\begin{loi_giai_cau}
\textbf{a)}
\[A = \sqrt{81} - \sqrt{16} = 9 - 4 = \boxed{5}.\]
\[B = |\sqrt{11}+2| - \sqrt{11} = \sqrt{11}+2-\sqrt{11} = \boxed{2}.\]

\textbf{b)}
\textbf{1.} Với $a > 0, a \neq 1$:
\[P = \left(\frac{\sqrt{a}}{\sqrt{a}(\sqrt{a}-1)} + \frac{2}{\sqrt{a}(\sqrt{a}-1)}\right) \cdot \frac{(\sqrt{a}-1)(\sqrt{a}+1)}{\sqrt{a}+2} = \frac{\sqrt{a}+2}{\sqrt{a}(\sqrt{a}-1)} \cdot \frac{(\sqrt{a}-1)(\sqrt{a}+1)}{\sqrt{a}+2} = \boxed{\frac{\sqrt{a}+1}{\sqrt{a}}}.\]

\textbf{2.} Khi $a = 3 + 2\sqrt{2} = (\sqrt{2}+1)^2$:
\[\sqrt{a} = \sqrt{(\sqrt{2}+1)^2} = \sqrt{2}+1.\]
Thay vào biểu thức $P$:
\[P = \frac{(\sqrt{2}+1)+1}{\sqrt{2}+1} = \frac{\sqrt{2}+2}{\sqrt{2}+1} = \frac{\sqrt{2}(\sqrt{2}+1)}{\sqrt{2}+1} = \boxed{\sqrt{2}}.\]
\end{loi_giai_cau}

\begin{de_cau}{Câu 2}{2,0 điểm}
\begin{enumerate}[label=\alph*)]
\item Giải phương trình: $x^2 - 3x - 10 = 0$.
\item Giải hệ phương trình: 
\[ \begin{cases} 2x + 3y = 8 \\ x - 2y = -3 \end{cases} \]
\end{enumerate}
\end{de_cau}

\begin{loi_giai_cau}
\textbf{a)} Phương trình $x^2 - 3x - 10 = 0$:
$\Delta = (-3)^2 - 4(1)(-10) = 9 + 40 = 49 = 7^2 > 0$.
Phương trình có hai nghiệm phân biệt:
\[x_1 = \frac{3+7}{2} = 5;\qquad x_2 = \frac{3-7}{2} = -2.\]
Vậy tập nghiệm là $\boxed{S = \{-2; 5\}}$.

\textbf{b)} Từ phương trình thứ hai: $x = 2y - 3$. Thế vào phương trình thứ nhất:
\[2(2y - 3) + 3y = 8 \iff 4y - 6 + 3y = 8 \iff 7y = 14 \iff y = 2.\]
Thay $y = 2$ vào $x = 2(2) - 3 = 1$.
Vậy hệ có nghiệm duy nhất $\boxed{(x; y) = (1; 2)}$.
\end{loi_giai_cau}

\begin{de_cau}{Câu 3}{1,5 điểm}
Cho phương trình bậc hai: $x^2 - 2(m+1)x + 2m = 0 \quad (1)$ ($m$ là tham số).
\begin{enumerate}[label=\alph*)]
\item Chứng minh rằng phương trình (1) luôn có hai nghiệm phân biệt với mọi giá trị của $m$.
\item Tìm $m$ để hai nghiệm $x_1, x_2$ thỏa mãn: $x_1^2 + x_2^2 = 12$.
\end{enumerate}
\end{de_cau}

\begin{loi_giai_cau}
\textbf{a)} Ta có:
\[\Delta' = [-(m+1)]^2 - 1 \cdot (2m) = m^2 + 2m + 1 - 2m = m^2 + 1.\]
Vì $m^2 \ge 0 \implies \Delta' = m^2 + 1 \ge 1 > 0$ với mọi $m$.
Vậy phương trình (1) luôn có hai nghiệm phân biệt $x_1, x_2$ với mọi $m$.

\textbf{b)} Theo định lý Vi-ét:
\[\begin{cases} x_1 + x_2 = 2(m+1) \\ x_1 x_2 = 2m \end{cases}\]
Ta có:
\[x_1^2 + x_2^2 = (x_1+x_2)^2 - 2x_1 x_2 = [2(m+1)]^2 - 2(2m) = 4(m^2+2m+1) - 4m = 4m^2 + 4m + 4.\]
Theo giả thiết:
\[4m^2 + 4m + 4 = 12 \iff 4m^2 + 4m - 8 = 0 \iff m^2 + m - 2 = 0 \iff (m-1)(m+2) = 0.\]
Suy ra $\boxed{m = 1}$ hoặc $\boxed{m = -2}$.
\end{loi_giai_cau}

\begin{de_cau}{Câu 5}{0,5 điểm}
Cho các số thực dương $a, b, c$ thỏa mãn $ab + bc + ca = 3$. Tìm giá trị nhỏ nhất của biểu thức:
\[P = \dfrac{1}{a^2+1} + \dfrac{1}{b^2+1} + \dfrac{1}{c^2+1}.\]
\end{de_cau}

\begin{loi_giai_cau}
Áp dụng bất đẳng thức Cauchy ngược dấu hoặc đánh giá:
Với mọi số dương $x$, ta có:
\[\frac{1}{x^2+1} = 1 - \frac{x^2}{x^2+1} \ge 1 - \frac{x^2}{2x} = 1 - \frac{x}{2}.\]
Áp dụng cho $a, b, c$:
\[P = \sum \frac{1}{a^2+1} \ge 3 - \frac{a+b+c}{2}.\]
Mặt khác, $(a+b+c)^2 \ge 3(ab+bc+ca) = 3 \times 3 = 9 \implies a+b+c \ge 3$.
Để tìm min, sử dụng bất đẳng thức quen thuộc:
Với $ab+bc+ca = 3$, ta có $\sum \dfrac{1}{a^2+1} \ge \dfrac{3}{2}$.
Thật vậy, khi $a = b = c = 1$:
\[P = \frac{1}{1^2+1} + \frac{1}{1^2+1} + \frac{1}{1^2+1} = \frac{1}{2} + \frac{1}{2} + \frac{1}{2} = \frac{3}{2}.\]
Chứng minh:
\[\frac{1}{a^2+1} - \frac{1}{2} = \frac{1-a^2}{2(a^2+1)} = \frac{1-a^2}{2(a^2+1)} \ge \frac{1-a}{2} \dots\]
Suy ra giá trị nhỏ nhất của $P$ là $\boxed{\dfrac{3}{2}}$, đạt được khi $a = b = c = 1$.
\end{loi_giai_cau}
